\documentclass{article}
\usepackage[utf8]{inputenc}
\usepackage[spanish]{babel}
\usepackage[letterpaper,top=2cm,bottom=2cm,left=3cm,right=3cm,marginparwidth=1.75cm]{geometry}
\usepackage{float}

\usepackage{amsmath}
\usepackage{amssymb}
\usepackage{graphicx}


\title{Top epigramas}
\author{
  Ramírez Ríos, Víctor Hugo\\
  \texttt{a221205528@unison.mx}
}
\date{Septiembre 2023}


\begin{document}

\maketitle

\section{10 que más me gustaron}
\begin{enumerate}
    \item Symmetry is a complexity-reducing concept (co-routines include subroutines); seek it everywhere.
    \par Entiendo por simetría a las partes de programas que suelen repetirse y hacer bulto en el código. Muchas veces se puede simplificar usando un mismo código en diferentes partes del programa. 
    \par Me gusta porque explica que no todo tiene que verse complejo para ser bueno, y usa la simetría como método para reducir la complejidad.
    \item A programming language is low level when its programs require attention to the irrelevant
    \par Aquí hace uso de bajo nivel como un concepto diferente al que conocemos, y solo llama un lenguaje bajo nivel cuando el lenguaje hace que desperdiciemos recursos en detalles irrelevantes.
    \par Me gusta porque no llama a lenguajes de tipo bajo nivel, sino que los clasifica más según el uso que les da el programador
    \item Recursion is the root of computation since it trades description for time
    \par Describe tanto la computación como la recursión como un concepto que cambia la simplicidad, por tiempo y eficiencia. 
    \par Me gusta porque describe la computación como algo simple, pero a cambio se sacrifican otros recursos.
    \item Every program has (at least) two purposes: the one for which it was written, and another for which it wasn’t.
    \par Describe como es que se reutilizan los programas. 
    \par Me gusta por que en el mundo de la programación, es muy común que se reutilizan códigos tanto para hacer modificaciones de programas o para hacer programas con propósitos diferentes. Es la idea central del código libre.
    \item  A program without a loop and a structured variable isn’t worth writing.
    \par Explica como es que los programas buenos generalmente son aquellos que requieran de cierta complejidad.
    \par Me gusta por que da una idea tangible sobre que vale la pena tomarse el tiempo para escribir. También, dice como es que puedes escribir un programa complejo pero si lo escribes sin estás dos cosas está mal planteado o ejecutado.
    \item There are two ways to write error-free programs; only the third one works.
    \par Hace referencia a como siempre hay errores en los programas.
    \par Me gusta por que aunque un programa marque que no tiene errores siempre hay posibilidad de mejorar.
    \item Fools ignore complexity. Pragmatists suffer it. Some can avoid it. Geniuses remove it.
    \par Habla sobre la complejidad en los programas, y como hay diferentes de manera de verla según el programador que seas.
    \par Me gusta porque una vez más le da importancia a la simplicidad en los programas. A su vez, menciona como es que siempre conviene eliminar esta complejidad.
    \item The computer is the ultimate polluter: its feces are indistinguishable from the food it produces. 
    \par Describe la salida y entrada de los programas, y como es que no es posible distinguir entre las que son correctas.
    \par Lo elegí por que no es muy claro en que se refiere a heces y comida.
    \item In programming, as in everything else, to be in error is to be reborn.
    \par Describe como equivocarse y tener errores es necesario para aprender y mejorar.
    \par Lo elegí por que al programar aprendemos más haciendo un programa, y corregir los errores, que copiando un programa y que salgo sin errores.
    \item If a listener nods his head when you’re explaining your program, wake him up. 
    \par Describe como los programas tienen un nivel de complejidad que requiere de cierto nivel de atención.
    \par Lo elegí por que siempre deben salir dudas a la hora que te muestren un programa, ya que es difícil entender un programa sin hacer preguntas.
\end{enumerate}

\section{10 que menos me gustaron}
\begin{enumerate}
    \item Syntactic sugar causes cancer of the semicolon.
    \par Menciona como es que en los lenguajes de programación no debería de haber elementos sintácticos que sean innecesarios.
    \par Lo elegí por que pienso que en algunos lenguajes esos elementos sintácticos denotan reglas muy definidas de como se escribe un programa.
    \item Optimization hinders evolution.
    \par Menciona como es que la evolución siempre conlleva a hacer programas peores programas.
    \par Lo elegí por que pienso que con el tiempo puede que hayan lenguajes más lentos, pero estos lenguajes nos ahorran tiempo y recursos para hacer programas más complejos.
    \item Once you understand how to write a program get someone else to write it.
    \par Menciona como es que solo es necesario entender a escribir una vez un programa para que sea un buen uso de tiempo. 
    \par Lo elegí por que pienso que saber escribir un programa, no significa que se entiendan todos los conceptos del mismo, y puedes volver a escribir uno similar y aprender y reforzar conocimientos.
    \item The eleventh commandment was ”Thou Shalt Compute” or ”Thou Shalt Not Compute” - I forget which.
    \par Menciona como es que hacer computaciones es una actividad que trae tanto cosas buenas como malas.
    \par Lo elegí por que no entiendo el objetivo de no saber si hacer computaciones es algo bueno o malo, si lo haces es por que te gusta.
    \item Some programming languages manage to absorb change, but withstand progress
    \par Menciona como es que hay algunos lenguajes de programación que a pesar de ser utilizados hoy en día tienen características que alentan el progreso.
    \par Lo elegí por que pienso que los lenguajes de programación de hoy en día se utilizan por que son los más convenientes en lo que hacen, y no por tener lealtad a ellos. Y aquellos que lo hacen por lealtad por lo general saben hacer muchos "trucos" con el lenguaje. 
    \item Giving up on assembly language was the apple in our Garden of Eden: Languages whose use squanders machine cycles are sinful. The LISP machine now permits LISP programmers to abandon bra and fig-leaf.
    \par Menciona como es que los lenguajes de programación modernos generan problemas que no eran recurrentes en los primordiales. 
    \par Lo elegí por que hace referencia a lenguajes que no conozco, y que tienen poca relevancia el día de hoy.
    
    \item Bringing computers into the home won’t change either one, but may revitalize the corner saloon.
    \par Entiendo que  estos epigramas son de hace tiempo, entonces pienso que este hace referencia a antes de que las laptops fueran una normalidad, y menciona como es que las computadores en casa serían más un entretenimiento que una zona de trabajo.
    \par Lo elegí por que pienso que cada quien usa su computadora para el uso que le vaya a dar, y tener una en casa trae más beneficios que consecuencias.
    
    \item If your computer speaks English, it was probably made in Japan.
    \par Pienso que este hace referencia a que cuando programamos lo hacemos en un lenguaje que no es el nuestro, y que muchas veces los países que no son ingleses programamos en inglés por que es la norma.
    \par Lo elegí por que no entiendo si hace referencia a la fuerza electrónica que tenía Japón hace años o a la significado que yo le di. Sí es el segundo, no encuentro un problema que normalicemos el uso de un lenguaje para la programación 
    \item Adapting old programs to fit new machines usually means adapting new machines to behave like old ones. 
    \par Lo entiendo como que los programas de antiguos nunca serán compatibles con las computadores del presente, y si queremos un balance entre ambas, una de los dos será perjudicada.
    \par Lo elegí por que hay programas viejos que a pesar que no han tenido cambios significantes hace años, siguen siendo la norma y la mejor opción a programas que son más nuevos.
    \item Computer Science is embarrassed by the computer. 
    \par Este lo entiendo más dirigido hacia el hecho que en la ciencias tenemos un enfoque más grande a las matemáticas, teoría, algoritmos, números, entre otras, que en sí enfocarnos en la computadora.
    \par Lo elegí por que pienso que muchas des esto que se ve en la ciencia es importante para entender el funcionamiento de las computadoras, y si no es así, al menos en mejorar nuestro pensamiento lógico y crítico, que es muy útil a la hora de entender como funciona una máquina. En realidad no pienso que sea incorrecto el epigrama, pero pienso que no es algo malo.

\end{enumerate}

\end{document}
